% =====================================================================
%  thmsmith Stage -1 proposal skeleton.
%  Written automatically by the thmsmith TS pipeline before Stage -1 runs.
%  Locked parameters: qid=eid\_dml\_orthogonality\_frontier, specialization=v1
%
%  HOW TO USE THIS FILE
%  --------------------
%  - The preamble (everything above the `% --- BEGIN CONTENT ---` marker)
%    and the `\section{...}` headers below are fixed by the pipeline.
%    Do NOT rewrite or rename them; the Step 3 reviewer subagent and the
%    downstream Stage 0 / Stage 1 prompts grep for these section titles.
%  - Each section contains exactly one `% TODO(stage-1 ...): ...`
%    placeholder. Replace each TODO with the content described for that
%    section in `CausalSmith/tools/src/discovery/prompts/stage_neg1_2_draft.txt`
%    ("REQUIRED OUTPUT FORMAT (Stage -1 proposal .tex)").
%  - The `\paragraph{Locked parameters.}` block in the metadata block
%    must pin the chosen `(qid, specialization, p, assignment, target,
%    regression)` precisely. If the proposal maps to the Q1 pool-mode,
%    reproduce that block verbatim from
%    `CausalSmith/doc/panel_formula_question_generator_note_with_common_prompts.tex`
%    (Q2..Q6 templates have been retired). Otherwise (free-form
%    `(qid, specialization)` — the default), define each parameter
%    explicitly here (no implicit conventions). Stage 0 quotes this block;
%    do not rephrase it after locking.
%  - Removing a section header, renaming a macro, or rewriting the
%    preamble is a regression: the file must still match this skeleton
%    after you finish.
% =====================================================================

\documentclass[11pt]{article}

\usepackage{amsmath,amssymb,amsthm,mathtools}
\usepackage{enumitem}
\usepackage[colorlinks=true,linkcolor=blue,citecolor=blue,urlcolor=blue]{hyperref}
\usepackage{natbib}

\newcommand{\E}{\mathbb{E}}
\renewcommand{\Pr}{\mathbb{P}}
\newcommand{\one}{\mathbf{1}}
\newcommand{\transpose}{\mathsf{T}}
\newcommand{\calH}{\mathcal{H}}
\newcommand{\calF}{\mathcal{F}}
\newcommand{\calR}{\mathcal{R}}
\newcommand{\R}{\mathbb{R}}
\newcommand{\plim}{\operatorname*{plim}}

\theoremstyle{definition}
\newtheorem{definition}{Definition}
\newtheorem{assumption}{Assumption}
\newtheorem{example}{Example}

\theoremstyle{plain}
\newtheorem{theorem}{Theorem}
\newtheorem{conjecture}{Conjecture}
\newtheorem{proposition}{Proposition}
\newtheorem{lemma}{Lemma}
\newtheorem{corollary}{Corollary}

\theoremstyle{remark}
\newtheorem{remark}{Remark}

\title{thmsmith Stage -1 proposal: \texttt{eid\_dml\_orthogonality\_frontier} / \texttt{v1}%
  \\ \large\textit{Primitive two-channel frontier for first-order HTE DML inference}}
\author{thmsmith Stage -1 proposer}
\date{\today}

\begin{document}
\maketitle

% --- BEGIN CONTENT ---  (everything above is fixed by the pipeline)

\paragraph{Locked parameters.}
This proposal uses the ExactID/free-form convention
\[
(\texttt{qid},\texttt{specialization},\text{SWIG/DAG},\text{treatment},
\text{target estimand},\text{identifying assumptions}).
\]
The locked tuple is
\[
\begin{gathered}
\texttt{qid}=\texttt{eid\_dml\_orthogonality\_frontier},\qquad
\texttt{specialization}=\texttt{v1},\\
\text{SWIG/DAG}=X\to D\to Y,\; X\to Y,\quad
\text{treatment }D\in\{0,1\},\\
\text{target estimand } \theta_0(\ell)=\E[\ell(X)\{Y(1)-Y(0)\}]
\text{ for }\ell\in\mathcal{L}_n,\\
\text{identifying assumptions}=\text{consistency, conditional unconfoundedness }
(Y(1),Y(0))\perp D\mid X,\\
\text{overlap } \underline e\le p_0(X)\le 1-\underline e,\text{ the Holder--entropy nuisance classes,}\\
\text{and the primitive Holder--entropy experiment in Assumptions
\ref{ass:unconfounded}--\ref{ass:primitive-geometry}.}
\end{gathered}
\]
Here \(p_0(x)=\Pr(D=1\mid X=x)\), \(g_{d,0}(x)=\E[Y\mid D=d,X=x]\), and
\(\mathcal{L}_n\) is a deterministic class of HTE report weights with entropy
integral \(\mathcal{E}_n\).  The focal mathematical object is the
\emph{two-channel orthogonality frontier}
\[
\mathfrak{F}_n(P)=
\max\{\mathfrak{D}_n(P),\mathfrak{S}_n(P)\},\qquad
\mathfrak{D}_n(P)=\sqrt n\,\delta_{g,n}(P)\delta_{p,n}(P),\quad
\mathfrak{S}_n(P)=\{\delta_{g,n}(P)+\delta_{p,n}(P)\}\mathcal{E}_n,
\]
where \(\delta_{g,n}\) and \(\delta_{p,n}\) are the \(L^2(P)\) nuisance radii
attainable by the learners over the locked Holder/entropy classes.  The
entropy factor is deliberately attached to the nuisance empirical-process
channel \(\mathfrak{S}_n\), not to the bounded population bilinear drift
\(\mathfrak{D}_n\).

\begin{abstract}
This proposal asks for the sharp boundary at which first-order Neyman
orthogonality stops delivering uniform Gaussian inference for heterogeneous
treatment-effect DML reports.  The target is not a new adjustment formula:
under unconfoundedness the weighted CATE reports are already identified by the
usual AIPW signal.  Version 5 splits the earlier iff claim into a conditional
bookkeeping theorem, where the upper and lower frontier channels are explicitly
assumed, and a flagship primitive conjecture, where those channels must be
derived from Holder smoothness, report entropy, and minimax lower-bound
geometry.  The proposed kernel is therefore not the conditional theorem but the
claim that the two-channel scalar \(\mathfrak{F}_n(P)\) is derivable as a
necessary-and-sufficient frontier over the primitive Holder--entropy
experiment.  Resolving it would tell DML users when first-order
orthogonalization is enough and when higher-order debiasing or a weaker report
class is mathematically necessary.
\end{abstract}

\section{Introduction}
\paragraph{Why the problem is important.}
First-order orthogonal DML is now a default route for treatment-effect
inference with machine-learned nuisances.  The standard message is robust but
qualitative: if nuisance products are small enough and the score class is tame
enough, the estimator behaves as if the nuisances were known.  For HTE reports,
that message leaves a real design question unanswered.  A confidence band, a
large collection of subgroup contrasts, or a rich CATE projection adds a second
nuisance channel: the population plug-in drift remains a bounded-envelope
product term, while the nuisance empirical process can be amplified by the
entropy of the report class.

\paragraph{The main finding (proposed).}
Theorem \ref{thm:conditional-frontier} records the conditional implication that
was previously overloaded as the flagship statement: if a high-level active
frontier bridge supplies both the sharp empirical-process upper bound and the
matching minimax lower-bound alternatives, then the iff boundary follows.
Conjecture \ref{conj:primitive-frontier} is the revised kernel.  It proposes
that the bridge can be proved, rather than assumed, for the primitive
Holder--entropy experiment: if
\(\mathfrak{F}_n(P)=\max\{\sqrt n\,\delta_{g,n}(P)\delta_{p,n}(P),
(\delta_{g,n}(P)+\delta_{p,n}(P))\mathcal{E}_n\}\) vanishes uniformly, the
\(\mathcal{L}_n\)-indexed DML process is oracle Gaussian; if it does not, a
minimax testing-distance criterion stays bounded away from zero through either
population bilinear drift or nuisance empirical-process leakage.

\paragraph{What is new.}
The closest published DML result, \citet[\S3]{ChernozhukovEtAl2018DML}, gives
sufficient product-rate and entropy conditions for orthogonal scores; it does
not claim necessity and does not identify a single operational scalar whose
crossing marks uniform failure.  \citet[\S2--\S3]{FosterSyrgkanis2023OSL}
show that orthogonality makes excess risk second order under entropy control,
but their object is risk rather than Gaussian inference.  \citet[\S3]{SemenovaChernozhukov2021DMLCATE}
give simultaneous inference for CATE summaries under high-level nuisance
conditions, while \citet[\S2--\S4]{KennedyBalakrishnanRobinsWasserman2024MinimaxCATE}
derive minimax CATE elbows that require higher-order ideas in low-smoothness
regimes.  \citet[abstract and pp.~1141--1142]{ArmstrongKolesar2021FiniteSampleATE}
and \citet[abstract]{MackeySyrgkanisZadik2018OML} are limitations comparators
that inform the conjecture rather than extension targets: the former ties low
smoothness and overlap to impossible \(\sqrt n\)-inference, and the latter
shows higher-order orthogonality can relax first-stage rate requirements.  The
epsilon-gap is now precise: Conjecture \ref{conj:primitive-frontier} must
replace the high-level frontier bridge in Assumption \ref{ass:frontier-active}
with primitive empirical-process and minimax lower-bound arguments, yielding an
iff frontier with the named two-channel object \(\mathfrak{F}_n(P)\).

\paragraph{How it positions in the literature.}
The anchor cluster is ExactID semiparametric inference under unconfoundedness.
The proposal consumes the exact-identification formula
\(\theta_0(\ell)=\E[\ell(X)\{g_{1,0}(X)-g_{0,0}(X)\}]\) and studies the
uniform inference consequences of estimating \(g_0\) and \(p_0\).  It sits at
the methodological-statistics edge of causal ML: DML, orthogonal statistical
learning, ADML/Riesz methods, minimax CATE theory, and honest finite-sample ATE
inference.

\paragraph{Implications for the community.}
The concrete downstream consumer is the simultaneous CATE and structural-function
workflow of \citet{SemenovaChernozhukov2021DMLCATE}: the proposed frontier would
say when its orthogonal-signal reports can be expanded to richer entropy-indexed
HTE classes without changing the debiasing order, and when the report class or
the score order must be changed.

\paragraph{How research benefits from the conclusion.}
If the conclusion resolves favorably, empirical DML papers should report or
bound \(\mathfrak{F}_n(P)\), not only separate nuisance rates, before claiming
uniform HTE confidence bands.

\section{Related work}
\citet{ChernozhukovEtAl2018DML} is the canonical DML anchor: it introduces
Neyman-orthogonal scores and cross-fitting, with sufficient product-rate and
entropy assumptions rather than a sharp uniform breakdown frontier.
\citet{SemenovaChernozhukov2021DMLCATE} moves DML toward HTE and CATE summaries
through orthogonal signals, best linear predictors, group averages, and
simultaneous inference; it is the closest positive HTE-inference comparator.
\citet{FosterSyrgkanis2023OSL} gives orthogonal statistical learning excess-risk
bounds under metric entropy and local complexity; it motivates the entropy
factor but does not deliver Gaussian approximation or lower bounds for DML
reports.
\citet{ChernozhukovChetverikovKato2014Gaussian} provides the Gaussian
approximation theorem planned for the oracle report process; this proposal uses
it as a proof-route input, not as prior art on orthogonal DML breakdown.
\citet{KennedyBalakrishnanRobinsWasserman2024MinimaxCATE} gives minimax CATE
rates and nuisance-smoothness elbows using higher-order influence-function
ideas; it is the main evidence that first-order orthogonality should have a
necessity boundary.
\citet[abstract and pp.~1141--1142]{ArmstrongKolesar2021FiniteSampleATE}
studies honest finite-sample ATE inference under smoothness, shape, and overlap
restrictions and states that feasible CIs can remain valid even when
\(\sqrt n\)-inference is impossible because of low smoothness or lack of
overlap; it supplies the necessity-oriented econometric precedent for not
treating DML rate conditions as merely technical.
\citet[abstract]{MackeySyrgkanisZadik2018OML} separates the power and limits of
first-order orthogonality by showing how higher-order orthogonality improves the
nuisance-rate requirement; it is the closest published limitations comparator
for the proposed first-order breakdown side.
\citet{ChernozhukovNeweySingh2022AutoDML} develops automatic DML and learned
Riesz representers for causal and structural effects; it is relevant because
its inference conditions depend on the learnability of a bias-correction object.
\citet{ChernozhukovNeweySingh2022Riesz} treats global and local linear
functionals, including nonregular local objects with growing Riesz norms; this
proposal uses the same regular-to-nonregular warning in an HTE report-class
frontier.
\citet{WagerAthey2018CausalForests}, \citet{AtheyTibshiraniWager2019GRF}, and
\citet{OprescuSyrgkanisWu2019ORF} provide forest-based HTE inference under
honesty, local estimating equations, and orthogonality; they are adjacent
algorithmic comparators rather than sharp Holder/entropy iff frontiers.
\citet{FarrellLiangMisra2021DNNInference} links deep neural-network rates to
semiparametric causal inference; it motivates modern nuisance classes where
approximation, entropy, and inference interact.

The precise published-axis gap for the v5 kernel is the absence of a primitive
lower-and-upper theorem.  The DML, CATE, OSL, and ADML papers above either give
sufficient high-level inference conditions, risk bounds, or minimax estimation
elbows; none derives both the empirical-process upper bound and the matching
minimax lower-bound alternatives for the two-channel scalar
\(\mathfrak{F}_n(P)\) over the same Holder--entropy HTE report experiment.

In-repo prior art is nonduplicative.  The auto-generated \texttt{doc/API.md}
catalogue is dominated by Panel Q1 basis-cleanliness theorems and records no
existing ExactID theorem for Holder/entropy DML frontiers.  The closest banked
CausalSmith predecessor is \texttt{eid\_dml\_orth\_uniform/v2}, whose rejected
kernel proposed a DML/LRO/ADML studentization-report frontier; reviewer signals
asked for explicit variance-report summands and scale-free spectral
normalization.  The present proposal changes the object from estimator-report
equivalence to a Holder/entropy first-order orthogonality breakdown frontier
for HTE reports, anchored by the named object \(\mathfrak{F}_n(P)\).

\section{Formal setup}
Let \(W=(Y,D,X)\sim P\), \(D\in\{0,1\}\), \(X\in[0,1]^d\), and
\(|Y|\le M\).  The SWIG/DAG is \(X\to D\to Y\) and \(X\to Y\).  Potential
outcomes \(Y(1),Y(0)\) satisfy consistency and conditional unconfoundedness.
Write
\[
g_{d,0}(x)=\E[Y\mid D=d,X=x],\qquad p_0(x)=\Pr(D=1\mid X=x),\qquad
\tau_0(x)=g_{1,0}(x)-g_{0,0}(x).
\]
For each report weight \(\ell\in\mathcal{L}_n\), the HTE report is
\[
\theta_0(\ell)=\E[\ell(X)\tau_0(X)].
\]
The identified AIPW signal is
\[
\Gamma(W;\eta_0)=g_{1,0}(X)-g_{0,0}(X)
+\frac{D}{p_0(X)}\{Y-g_{1,0}(X)\}
-\frac{1-D}{1-p_0(X)}\{Y-g_{0,0}(X)\},
\]
where \(\eta_0=(g_{1,0},g_{0,0},p_0)\).  Then
\(\theta_0(\ell)=\E[\ell(X)\Gamma(W;\eta_0)]\).  The first-order DML estimator
uses cross-fitted nuisance estimates \(\hat\eta=(\hat g_1,\hat g_0,\hat p)\):
\[
\hat\theta_n(\ell)=\mathbb{P}_n[\ell(X)\Gamma(W;\hat\eta_{-k(i)})],
\qquad \ell\in\mathcal{L}_n.
\]
The nuisance empirical-process remainder is
\[
\Delta^{\mathrm{ep}}_n(\ell)
=
(\mathbb{P}_n-P)\!\left[
\ell(X)\{\Gamma(W;\hat\eta_{-k(i)})-\Gamma(W;\eta_0)\}
\right],
\qquad \ell\in\mathcal{L}_n,
\]
where the fold index emphasizes that the nuisance estimate in the score for
observation \(i\) is trained off the fold containing \(i\).
Let \(\delta_{g,n}(P)=\|\hat g_1-g_{1,0}\|_{2,P}+\|\hat g_0-g_{0,0}\|_{2,P}\)
and \(\delta_{p,n}(P)=\|\hat p-p_0\|_{2,P}\), interpreted as deterministic
upper envelopes over the locked learner/nuisance class.  For Holder balls
\(\mathcal{H}(\alpha,L,d)\), the canonical minimax radii are
\[
r_g=\frac{\alpha_g}{2\alpha_g+d},\quad
r_p=\frac{\alpha_p}{2\alpha_p+d},\quad
\delta_{g,n}\asymp n^{-r_g},\quad \delta_{p,n}\asymp n^{-r_p},
\]
up to logarithmic factors.  The HTE report class has entropy integral
\[
\mathcal{E}_n=
1+\int_0^1\sqrt{\log N(\varepsilon\|\mathcal{L}_n\|_{2,P},
\mathcal{L}_n,L^2(P))}\,d\varepsilon,
\]
uniformly over \(P\in\mathcal{P}_n\).  The focal frontier is
\[
\mathfrak{F}_n(P)=\max\{\mathfrak{D}_n(P),\mathfrak{S}_n(P)\},\qquad
\mathfrak{D}_n(P)=\sqrt n\,\delta_{g,n}(P)\delta_{p,n}(P),\quad
\mathfrak{S}_n(P)=\{\delta_{g,n}(P)+\delta_{p,n}(P)\}\mathcal{E}_n.
\]
When \(\mathcal{E}_n\asymp n^\zeta\) and the nuisance radii have the Holder
rates above, the exponent version of the boundary is the pair of inequalities
\[
r_g+r_p>\frac12,\qquad \min\{r_g,r_p\}>\zeta.
\]
Let
\[
\mathbb{G}_n(\ell)=\sqrt n\,(\mathbb{P}_n-P)[\ell(X)\{\Gamma(W;\eta_0)-\theta_0(\ell)\}]
\]
be the oracle Gaussian target process.  For a pair of nuisance perturbations
\(u=(u_1,u_0)\) and \(v\), define the ideal-denominator second-order remainder
bilinear map
\[
R_P(\ell;u,v)=
\E_P\!\left[
\ell(X)v(X)
\left\{\frac{u_1(X)}{p_0(X)}+\frac{u_0(X)}{1-p_0(X)}\right\}
\right].
\]
\[
\widetilde R_P(\ell;u,v,\bar p)=
\E_P\!\left[
\ell(X)v(X)
\left\{\frac{u_1(X)}{\bar p(X)}+\frac{u_0(X)}{1-\bar p(X)}\right\}
\right]
\]
is the exact plug-in drift map for an estimated propensity \(\bar p\), and
\[
Q_P(\ell;u,v,\bar p)=\widetilde R_P(\ell;u,v,\bar p)-R_P(\ell;u,v)
\]
is the estimated-denominator residual.
The curvature constant \(c_R(P)\) in Assumption \ref{ass:primitive-geometry}
is the operator norm of \(R_P\) on the active Holder tangent cone.  The
empirical leakage constant \(c_{\mathrm{ep}}(P)\) is the corresponding
Sudakov/minoration constant for the nuisance empirical-process class in the
same primitive experiment.
The minimax oracle-centering defect used in Conjecture
\ref{conj:primitive-frontier} is
\[
\mathcal{K}_n(\mathcal{P}_n,\mathcal{L}_n)
=
\inf_{\hat\theta}
\sup_{P\in\mathcal{P}_n}
d_{\mathrm{BL}}\!\left(
\mathcal{L}_P\!\left(
\sup_{\ell\in\mathcal{L}_n}
\left|\sqrt n\{\hat\theta(\ell)-\theta_0(\ell)\}-\mathbb{G}_n(\ell)\right|
\right),\delta_0\right),
\]
where \(d_{\mathrm{BL}}\) is bounded-Lipschitz distance and \(\delta_0\) is
point mass at zero.  A positive liminf of \(\mathcal{K}_n\) is the formal
testing-distance criterion for uniform first-order Gaussian failure.

\begin{quote}
\begin{verbatim}
Locked tuple:
(qid, specialization, SWIG/DAG, treatment, target estimand, identifying assumptions)
qid = eid_dml_orthogonality_frontier
specialization = v1
SWIG/DAG = X -> D -> Y and X -> Y
treatment = binary D
target estimand = theta_0(ell)=E[ell(X){Y(1)-Y(0)}], ell in L_n
identifying assumptions = consistency, conditional unconfoundedness,
overlap, Holder--entropy nuisance classes, and Assumptions 1--7.
\end{verbatim}
\end{quote}

\section{Assumptions}
\begin{assumption}[Unconfounded exact-ID model]\label{ass:unconfounded}
For every \(P\in\mathcal{P}_n\), \(Y=DY(1)+(1-D)Y(0)\),
\((Y(1),Y(0))\perp D\mid X\), \(X\in[0,1]^d\), \(|Y|\le M\), and
\(\underline e\le p_0(X)\le1-\underline e\) almost surely for a fixed
\(\underline e>0\).
\end{assumption}

\begin{assumption}[Holder nuisance and learner radii]\label{ass:holder}
The outcome regressions \(g_{1,0},g_{0,0}\) belong to a Holder ball
\(\mathcal{H}(\alpha_g,L_g,d)\), the propensity \(p_0\) belongs to
\(\mathcal{H}(\alpha_p,L_p,d)\), and the cross-fitted learners have uniform
\(L^2(P)\) envelopes \(\delta_{g,n}(P)\) and \(\delta_{p,n}(P)\).  In the
canonical Holder-rate specialization,
\(\delta_{g,n}\asymp n^{-\alpha_g/(2\alpha_g+d)}\) and
\(\delta_{p,n}\asymp n^{-\alpha_p/(2\alpha_p+d)}\).
\end{assumption}

\begin{assumption}[Report-class entropy]\label{ass:entropy}
The deterministic HTE report class \(\mathcal{L}_n\) is uniformly bounded by
\(L_\ell\), centered when required by the covariance normalization, and has
finite entropy integral \(\mathcal{E}_n\) as defined in Section 3.  Polynomial
entropy classes with \(\log N(\varepsilon,\mathcal{L}_n,L^2(P))\lesssim
a_n\varepsilon^{-v_\ell}\) are allowed; \(\mathcal{E}_n\) records the resulting
growth.
\end{assumption}

\begin{assumption}[Cross-fitting and oracle-process regularity]\label{ass:crossfit}
The sample is split into a fixed number of folds, nuisance estimates used for
observation \(i\) are trained off fold \(k(i)\), and the oracle process
\(\mathbb{G}_n\) is the target oracle empirical process on
\(\mathcal{L}_n\).  The primitive oracle-process route is the
Chernozhukov--Chetverikov--Kato Gaussian approximation for suprema of empirical
processes, plus the corresponding maximal inequality for VC/polynomial-entropy
report classes.  This assumption fixes fold independence and the intended
oracle comparison theorem; it does not assume the sharp bound for
\(\Delta^{\mathrm{ep}}_n\) at scale \(\mathfrak{S}_n(P)\).
\end{assumption}

\begin{assumption}[Estimated-denominator residual]\label{ass:denominator}
The estimated propensity obeys the same fixed overlap bound as \(p_0\) with
probability approaching one, and the denominator-substitution residual is
negligible on the upper-bound side:
\[
\sqrt n\sup_{\ell\in\mathcal{L}_n}
\left|Q_P(\ell;\hat g-g_0,\hat p-p_0,\hat p)\right|=o_P(1).
\]
A sufficient primitive envelope is
\(\sqrt n\,\delta_{g,n}(P)\delta_{p,n}(P)^2=o(1)\), together with bounded
report weights and fixed overlap.
\end{assumption}

\begin{assumption}[Primitive Holder--entropy tangent geometry]\label{ass:primitive-geometry}
The Holder balls admit wavelet or spline tangent sets
\(\mathcal{T}_{g,n}\times\mathcal{T}_{p,n}\) with finite support dimension
growing at the minimax resolutions, and the report class \(\mathcal{L}_n\)
contains a matching finite packet subfamily with entropy integral
\(\mathcal{E}_n\).  Define
\[
c_R(P)=
\sup_{\ell\in\mathcal{L}_n,\; \|u\|_{2,P}\le1,\; \|v\|_{2,P}\le1}
|R_P(\ell;u,v)|
\]
restricted to the active tangent packet, and define \(c_{\mathrm{ep}}(P)\) as
the Sudakov/minoration constant for the centered class
\(\{\ell(X)(\Gamma(W;\eta_0+t h)-\Gamma(W;\eta_0)):
\ell\in\mathcal{L}_n,h\in\mathcal{T}_{g,n}\times\mathcal{T}_{p,n}\}\) at the
local nuisance scale.  On the active channel under consideration,
\(\inf_{P\in\mathcal{P}_n}\max\{c_R(P),c_{\mathrm{ep}}(P)\}\ge c_0>0\), where
\(c_0\) is a fixed structural constant independent of \(n\).
\end{assumption}

\begin{assumption}[Conditional active-frontier bridge]\label{ass:frontier-active}
There exist nuisance estimators and a tangent-cone collection such that
(i) the empirical-process remainder of the AIPW exact-identification score is
bounded above by \(\mathfrak{S}_n(P)\), up to logarithmic factors,
(ii) the oracle Gaussian approximation holds whenever
\(\mathfrak{F}_n(P)\to0\), and (iii) active local alternatives realize
matching minimax lower bounds with constants \(c_R,c_{\mathrm{ep}}>0\) at the
same scale \(\mathfrak{F}_n(P)\).  This is the consolidated
\texttt{A-frontier-active} bridge supplied after the Stage 0.5 rewind; it is
used only for the conditional theorem and is exactly what Conjecture
\ref{conj:primitive-frontier} proposes to derive from primitive arguments.
\end{assumption}

\section{Main results}
\begin{theorem}[Theorem 1: exact-ID score and exact plug-in remainder]\label{thm:score-remainder}
Under Assumptions \ref{ass:unconfounded}--\ref{ass:denominator}, for every
\(\ell\in\mathcal{L}_n\),
\(\theta_0(\ell)=\E[\ell(X)\Gamma(W;\eta_0)]\), and the cross-fit plug-in
expansion satisfies
\[
\hat\theta_n(\ell)-\theta_0(\ell)
=(\mathbb{P}_n-P)\ell(X)\{\Gamma(W;\eta_0)-\theta_0(\ell)\}
+\widetilde R_P(\ell;\hat g-g_0,\hat p-p_0,\hat p)
+\Delta^{\mathrm{ep}}_n(\ell)+o_P(n^{-1/2})
\]
uniformly in \(\ell\in\mathcal{L}_n\).  The displayed empirical-process summand
is intentionally left explicit: bounding it at the sharp scale
\(\mathfrak{S}_n(P)\) is not assumed in this theorem and is part of the
primitive frontier conjecture below.  Moreover, under Assumption
\ref{ass:denominator},
\[
\sup_{\ell\in\mathcal{L}_n}
\left|
\sqrt n\{\widetilde R_P(\ell;\hat g-g_0,\hat p-p_0,\hat p)
-R_P(\ell;\hat g-g_0,\hat p-p_0)\}
\right|=o_P(1),
\]
so the ideal-denominator bilinear map \(R_P\) is the valid frontier object.
\end{theorem}
\paragraph{Why this is new and worth studying.}
(a) This is not the flagship novelty: it records the exact-ID AIPW identity used
by \citet{ChernozhukovEtAl2018DML} and \citet{SemenovaChernozhukov2021DMLCATE},
but isolates the \(\mathcal{L}_n\)-indexed bilinear remainder \(R_P\) that those
papers control by sufficient conditions.  (b) The consumer is Conjecture
\ref{conj:primitive-frontier}: without this uniform remainder object, the proposed
frontier would be a rate slogan rather than a theorem-level target.

\begin{theorem}[Theorem 2: conditional active-frontier theorem]\label{thm:conditional-frontier}
Under Assumptions \ref{ass:unconfounded}--\ref{ass:frontier-active},
first-order orthogonal DML is uniformly asymptotically linear and Gaussian for
the \(\mathcal{L}_n\)-indexed HTE report process if and only if
\[
\sup_{P\in\mathcal{P}_n}\mathfrak{F}_n(P)
=\sup_{P\in\mathcal{P}_n}\max\{\sqrt n\,\delta_{g,n}(P)\delta_{p,n}(P),
\{\delta_{g,n}(P)+\delta_{p,n}(P)\}\mathcal{E}_n\}
\longrightarrow0 .
\]
The theorem is conditional: the upper empirical-process scale, the oracle
Gaussian approximation, and the matching lower-bound alternatives are supplied
by Assumption \ref{ass:frontier-active}.  It should be read as a theorem split
that makes the earlier high-level frontier bridge explicit, not as the
flagship novelty.
\end{theorem}
\paragraph{Why this is new and worth studying.}
(a) This is a correction, not the new flagship claim: the Stage 0.5 rejection
identified v4 Assumptions 4 and 6 as encoding the upper and lower sides of the
frontier.  Theorem
\ref{thm:conditional-frontier} demotes that content to a conditional
bookkeeping statement under the consolidated bridge Assumption
\ref{ass:frontier-active}.  (b) The consumer is Stage 0's derivation pipeline:
it can formalize the conditional theorem without mistaking it for the
unconditional contribution.

\begin{conjecture}[Conjecture 1: primitive two-channel minimax frontier (NOT YET PROVED)]\label{conj:primitive-frontier}\label{conj:frontier}
Under Assumptions \ref{ass:unconfounded}--\ref{ass:primitive-geometry}, and
without invoking Assumption \ref{ass:frontier-active}, the conclusions of the
active-frontier bridge can be derived from primitive empirical-process and
minimax arguments for the locked Holder--entropy experiment.  Equivalently,
first-order orthogonal DML is uniformly asymptotically linear and Gaussian for
the \(\mathcal{L}_n\)-indexed HTE report process if and only if
\[
\sup_{P\in\mathcal{P}_n}\mathfrak{F}_n(P)\to0.
\]
The lower-bound side is formalized by the testing-distance criterion
\(\mathcal{K}_n(\mathcal{P}_n,\mathcal{L}_n)\) from Section 3: if
\(\limsup_n\sup_{P\in\mathcal{P}_n}\mathfrak{F}_n(P)>0\), then
\[
\liminf_{n\to\infty}\mathcal{K}_n(\mathcal{P}_n,\mathcal{L}_n)>0
\]
for an open subset of the locked Holder nuisance class.  If
\(\mathfrak{D}_n\) is the active channel, the proof must construct local
least-favorable \(u_n,v_n\) with a nonzero bilinear translation \(R_P\).  If
\(\mathfrak{S}_n\) is the active channel, the proof must produce a
Sudakov/Haar empirical-leakage witness at scale
\((\delta_{g,n}+\delta_{p,n})\mathcal{E}_n\).  Faster oracle learners outside
the active error class are not claimed to fail.
\end{conjecture}
\paragraph{Why this is new and worth studying.}
(a) \citet[\S3]{ChernozhukovEtAl2018DML} gives sufficient product-rate
conditions, \citet[\S2--\S3]{FosterSyrgkanis2023OSL} gives excess-risk rather
than inference thresholds, and \citet[\S3]{SemenovaChernozhukov2021DMLCATE}
states high-level simultaneous-inference conditions; none states or proves an
iff frontier with the named two-channel object \(\mathfrak{F}_n(P)\), a
bounded-Lipschitz minimax failure criterion \(\mathcal{K}_n\), and open-class
lower-bound witnesses separating deterministic drift from empirical leakage.
(b) The consumer is the HTE simultaneous-inference workflow in
\citet{SemenovaChernozhukov2021DMLCATE}: resolving this conjecture tells that
workflow exactly when a richer report class invalidates first-order DML.

\begin{conjecture}[Conjecture 2: Holder exponent corollary of the primitive frontier (NOT YET PROVED)]\label{conj:holder}
In the canonical Holder-rate specialization of Assumption \ref{ass:holder} with
\(\mathcal{E}_n\asymp n^\zeta\), the frontier in Conjecture
\ref{conj:frontier} is equivalent, up to logarithmic factors, to
\[
\frac{\alpha_g}{2\alpha_g+d}+\frac{\alpha_p}{2\alpha_p+d}
>\frac12,\qquad
\min\left\{\frac{\alpha_g}{2\alpha_g+d},\frac{\alpha_p}{2\alpha_p+d}\right\}
>\zeta.
\]
Below this boundary, first-order DML cannot be uniformly Gaussian over the
locked Holder/entropy class; above it, first-order DML is uniformly Gaussian
provided Conjecture \ref{conj:primitive-frontier}'s primitive empirical-process
and oracle-Gaussian upper bound has been proved.
\end{conjecture}
\paragraph{Why this is new and worth studying.}
(a) \citet[\S2--\S4]{KennedyBalakrishnanRobinsWasserman2024MinimaxCATE}
identifies nuisance-smoothness elbows for minimax CATE estimation and uses
higher-order ideas, while
\citet[abstract and pp.~1141--1142]{ArmstrongKolesar2021FiniteSampleATE} shows
that low smoothness or lack of overlap can preclude \(\sqrt n\)-inference and
\citet[abstract]{MackeySyrgkanisZadik2018OML} shows higher-order orthogonality
can relax first-stage rate requirements; none gives this first-order DML
uniform-inference exponent boundary for HTE report classes.  This result is an
exponent translation of Conjecture \ref{conj:primitive-frontier}, not an
independent derivation of the primitive frontier.
(b) The consumer is the higher-order orthogonalization program in
\citet{KennedyBalakrishnanRobinsWasserman2024MinimaxCATE}: the conjecture
pinpoints the exponent region where first-order scores must be replaced or the
target report class weakened.

\section{Sanity-check corollaries}
\begin{corollary}[ATE singleton reduction]\label{cor:ate}
If \(\mathcal{L}_n=\{1\}\), then \(\mathcal{E}_n\asymp1\) and Conjecture
\ref{conj:frontier} reduces to the familiar DML product-rate condition
\(\sqrt n\,\delta_{g,n}\delta_{p,n}\to0\), matching
\citet{ChernozhukovEtAl2018DML} up to logarithmic factors.
\end{corollary}

\begin{corollary}[Fixed subgroup-report reduction]\label{cor:fixed-subgroups}
If \(\mathcal{L}_n\) is a fixed finite collection of subgroup indicators or
Haar cells, then \(\mathcal{E}_n=O(1)\).  Under the primitive frontier in
Conjecture \ref{conj:primitive-frontier}, the entropy channel is constant and
the boundary collapses to the ATE-style product-rate condition plus
\(\delta_{g,n}+\delta_{p,n}\to0\), matching the fixed-dimensional report
intuition in \citet{SemenovaChernozhukov2021DMLCATE}.  If the number of Haar
cells grows with \(n\), \(\mathcal{E}_n\) reappears and the empirical-leakage
channel becomes the binding constraint.
\end{corollary}

\begin{corollary}[Equal-smoothness Holder elbow]\label{cor:equal-holder}
If \(\alpha_g=\alpha_p=\alpha\) and \(\mathcal{E}_n\asymp1\), Conjecture
\ref{conj:holder} reduces to \(\alpha>d/2\).  If
\(\mathcal{E}_n\asymp n^\zeta\), it reduces to
\(2\alpha/(2\alpha+d)>1/2\) and \(\alpha/(2\alpha+d)>\zeta\), so richer HTE
reporting requires enough nuisance smoothness to control the empirical-process
channel in addition to the usual product-rate channel.
\end{corollary}

\begin{example}[Concrete non-corner witness class]\label{ex:witness}
Take \(X\sim \mathrm{Unif}[0,1]^d\), \(p_0(x)=1/2+\rho\prod_{j=1}^d
\sin(2\pi x_j)\) with \(0<\rho<1/4\), and outcome regressions equal to
finite wavelet packets at resolution \(J_n\) scaled to lie in the Holder balls.
Let \(\mathcal{L}_n\) be the span of the corresponding Haar cells with
\(\mathcal{E}_n\asymp n^\zeta\).  The bounded-envelope check is explicit:
\[
\sup_{\ell\in\mathcal{L}_n}|R_P(\ell;u,v)|
\le C(\underline e,L_\ell)\|u\|_{2,P}\|v\|_{2,P},
\]
so the witness only asserts \(R_P\ge c_R\delta_{g,n}\delta_{p,n}\), with no
entropy multiplier in the population drift.  Entropy enters through the Haar
maximal inequality
\[
\sqrt n\sup_{\ell\in\mathcal{L}_n}|\Delta^{\mathrm{ep}}_n(\ell)|
\gtrsim c_{\mathrm{ep}}\{\delta_{g,n}+\delta_{p,n}\}\mathcal{E}_n
\]
on the active finite-packet subfamily.  Away from a finite collection of
coefficient hyperplanes, either \(c_R(P)>0\) or \(c_{\mathrm{ep}}(P)>0\), giving
the proposed open-class witness for the necessity half of Conjecture
\ref{conj:frontier}.
\end{example}

\section{Proof-sketch route}
Theorem \ref{thm:score-remainder} is the standard AIPW identity plus an exact
plug-in drift calculation: conditional on \(X\), the drift equals
\((\hat p-p_0)\{(\hat g_1-g_{1,0})/\hat p+(\hat g_0-g_{0,0})/(1-\hat p)\}\),
and Assumption \ref{ass:denominator} replaces estimated denominators by
population denominators at \(o(n^{-1/2})\) scale.  Theorem
\ref{thm:conditional-frontier} then follows by substituting the consolidated
active-frontier bridge into that expansion; this is the conditional theorem
split requested by the Stage 0.5 review.  The sufficiency half of Conjecture
\ref{conj:primitive-frontier} must replace that bridge with a
bounded-envelope Cauchy--Schwarz bound for the \(\mathcal{L}_n\)-indexed
bilinear population remainder, a maximal inequality for
\(\Delta^{\mathrm{ep}}_n\) at scale \(\mathfrak{S}_n\), a high-dimensional
Gaussian comparison for the oracle process, and anti-concentration for the
multiplier statistic.  The necessity half should split by active channel.  If
\(\mathfrak{D}_n\) is active, the primitive tangent geometry in Assumption
\ref{ass:primitive-geometry} should be used to construct local least-favorable
Holder perturbations \(u_n,v_n\) whose bilinear remainder has size
\(\delta_{g,n}\delta_{p,n}\), followed by a two-point lower-bound argument for
\(\mathcal{K}_n\).  If \(\mathfrak{S}_n\) is active, the Haar/wavelet packet in
Example \ref{ex:witness} supplies the intended Sudakov or high-dimensional
maximal-inequality lower bound for the nuisance empirical process itself.  When
\(\mathfrak{F}_n\to\infty\), the centered process has a non-negligible
active-channel error; when \(\mathfrak{F}_n\to c\in(0,\infty)\), the active
channel produces a fixed nonzero departure from the oracle Gaussian process,
so the oracle-centered multiplier law misses the limit even though the failure
is not explosive.  The Holder exponent statement follows by
substituting the classical nonparametric \(L^2\) rates and the entropy-integral
growth of \(\mathcal{L}_n\).  Boundary cases needing explicit algebra are the
singleton ATE report, the finite-dimensional subgroup report, the Haar/wavelet
witness in Example \ref{ex:witness}, and weak-overlap sequences where
\(\underline e\) is allowed to drift.

\section{Honest scope flags}
Theorem \ref{thm:score-remainder} is intended as routine setup, not the
flagship claim.  Theorem \ref{thm:conditional-frontier} is also not the
flagship claim: it is a conditional theorem under Assumption
\ref{ass:frontier-active}, which explicitly assumes the upper and lower
frontier channels.  Conjectures \ref{conj:primitive-frontier} and
\ref{conj:holder} are not yet proved; the primitive maximal inequality,
Sudakov/Haar lower-bound witness, and two-point minimax testing-distance
argument are the open mathematical steps.  The proposal does not claim that
first-order DML is invalid in all low-smoothness problems: it claims invalidity
over the locked open Holder class when either the population curvature constant
or the empirical leakage constant is bounded away from zero.  The most delicate
primitive step is replacing Assumption \ref{ass:frontier-active}, because
Gaussian approximation for a growing HTE report class can impose restrictions
separate from the population orthogonality remainder.  Version 2 patches the exact
plug-in denominator algebra and narrows the lower-bound half to a
minimax/adversarial first-order experiment.  Version 3 makes the nuisance
empirical-process remainder explicit and assumes its uniform negligibility.
Version 4 withdraws the entropy-multiplied population-drift normalization and
places entropy only in the empirical-process channel; no kernel claims have
been withdrawn.  Version 5 demotes the earlier iff frontier to a conditional
theorem under Assumption \ref{ass:frontier-active} and makes the primitive
empirical-process/minimax derivation the open flagship conjecture.

\section{Iteration trail}
\begin{itemize}[leftmargin=*]
\item v1 (cold-start) -- initial draft of the entropy-remainder frontier for
first-order HTE DML; prior verdict: none.
\item v2 (revise) -- patched the exact plug-in remainder, lower-bound
quantifiers, Armstrong--Kolesar anchor, and intermediate-boundary proof route;
prior verdict: REVISE.
\item v3 (revise) -- displayed the nuisance empirical-process remainder in
Theorem \ref{thm:score-remainder} and added its uniform \(o_P(n^{-1/2})\)
condition to Assumption \ref{ass:crossfit}; prior verdict: REVISE.
\item v4 (revise) -- split the old entropy-multiplied drift into the
population channel \(\mathfrak{D}_n\) and the empirical-process channel
\(\mathfrak{S}_n\), removing entropy growth from the bounded bilinear
population remainder; prior verdict: REVISE.
\item v5 (revise) -- theorem-split after Stage 0.5: demoted the iff frontier
to a conditional active-frontier theorem and made the primitive
empirical-process/minimax derivation the open flagship conjecture; prior
verdict: REJECT.
\end{itemize}

\section*{Bibliography}
\begin{thebibliography}{99}
\bibitem[Armstrong and Kolesar(2021)]{ArmstrongKolesar2021FiniteSampleATE}
Armstrong, T. B. and Kolesar, M. (2021).
\newblock Finite-sample optimal estimation and inference on average treatment effects under unconfoundedness.
\newblock \emph{Econometrica} 89(3), 1141--1177.

\bibitem[Athey, Tibshirani, and Wager(2019)]{AtheyTibshiraniWager2019GRF}
Athey, S., Tibshirani, J., and Wager, S. (2019).
\newblock Generalized random forests.
\newblock \emph{Annals of Statistics} 47(2), 1148--1178.

\bibitem[Chernozhukov et~al.(2018)]{ChernozhukovEtAl2018DML}
Chernozhukov, V., Chetverikov, D., Demirer, M., Duflo, E., Hansen, C., Newey, W., and Robins, J. (2018).
\newblock Double/debiased machine learning for treatment and structural parameters.
\newblock \emph{Econometrics Journal} 21(1), C1--C68.

\bibitem[Chernozhukov, Chetverikov, and Kato(2014)]{ChernozhukovChetverikovKato2014Gaussian}
Chernozhukov, V., Chetverikov, D., and Kato, K. (2014).
\newblock Gaussian approximation of suprema of empirical processes.
\newblock \emph{Annals of Statistics} 42(4), 1564--1597.

\bibitem[Chernozhukov, Newey, and Singh(2022a)]{ChernozhukovNeweySingh2022AutoDML}
Chernozhukov, V., Newey, W., and Singh, R. (2022a).
\newblock Automatic debiased machine learning of causal and structural effects.
\newblock \emph{Econometrica} 90(3), 967--1027.

\bibitem[Chernozhukov, Newey, and Singh(2022b)]{ChernozhukovNeweySingh2022Riesz}
Chernozhukov, V., Newey, W., and Singh, R. (2022b).
\newblock Debiased machine learning of global and local parameters using regularized Riesz representers.
\newblock \emph{Econometrics Journal} 25(3), 576--601.

\bibitem[Farrell, Liang, and Misra(2021)]{FarrellLiangMisra2021DNNInference}
Farrell, M. H., Liang, T., and Misra, S. (2021).
\newblock Deep neural networks for estimation and inference.
\newblock \emph{Econometrica} 89(1), 181--213.

\bibitem[Foster and Syrgkanis(2023)]{FosterSyrgkanis2023OSL}
Foster, D. J. and Syrgkanis, V. (2023).
\newblock Orthogonal statistical learning.
\newblock \emph{Annals of Statistics} 51(3), 879--908.

\bibitem[Kennedy et~al.(2024)]{KennedyBalakrishnanRobinsWasserman2024MinimaxCATE}
Kennedy, E. H., Balakrishnan, S., Robins, J. M., and Wasserman, L. (2024).
\newblock Minimax rates for heterogeneous causal effect estimation.
\newblock \emph{Annals of Statistics} 52(2), 793--817.

\bibitem[Mackey, Syrgkanis, and Zadik(2018)]{MackeySyrgkanisZadik2018OML}
Mackey, L., Syrgkanis, V., and Zadik, I. (2018).
\newblock Orthogonal machine learning: Power and limitations.
\newblock \emph{Proceedings of the 35th International Conference on Machine Learning}, PMLR 80, 3375--3383.

\bibitem[Oprescu, Syrgkanis, and Wu(2019)]{OprescuSyrgkanisWu2019ORF}
Oprescu, M., Syrgkanis, V., and Wu, Z. S. (2019).
\newblock Orthogonal random forest for causal inference.
\newblock \emph{Proceedings of the 36th International Conference on Machine Learning}.

\bibitem[Semenova and Chernozhukov(2021)]{SemenovaChernozhukov2021DMLCATE}
Semenova, V. and Chernozhukov, V. (2021).
\newblock Debiased machine learning of conditional average treatment effects and other causal functions.
\newblock \emph{Econometrics Journal} 24(2), 264--289.

\bibitem[Wager and Athey(2018)]{WagerAthey2018CausalForests}
Wager, S. and Athey, S. (2018).
\newblock Estimation and inference of heterogeneous treatment effects using random forests.
\newblock \emph{Journal of the American Statistical Association} 113(523), 1228--1242.
\end{thebibliography}

\end{document}
